\BeleskeID{08-s059}
% element: 08-s059:e04
% element: 08-s059:notes
Impuls se rastavlja na prvi skok prema $V_{OH}$ i drugi, negativan skok amplitude $V_{OH}-V_{OL}$ u $t_2$. Drugi odziv pre trenutka $t_2$ je nula. Po početku tog skoka njegov opšti oblik je
\[u_{o2}(t)=-(V_{OH}-V_{OL})+\bigl(0-(-(V_{OH}-V_{OL}))\bigr)e^{-\frac{t-t_2}{\tau}},\quad t\geq t_2.\]
Sređivanjem se dobija izraz za $u_{o2}$ prikazan na slajdu.

Pre početka impulsa izlaz je $V_{OL}$, a dok drugi skok još nije nastupio, rezultat je samo prvi odziv:
\begin{align*}
u_o(t)&=V_{OL},\quad t<t_1,\\
u_o(t)&=V_{OH}+(V_{OL}-V_{OH})e^{-\frac{t-t_1}{\tau}},\quad t_1\leq t\leq t_2.
\end{align*}
Dva doprinosa pri $t\geq t_2$ imaju oblike
\[V_{OH}+(V_{OL}-V_{OH})e^{-\frac{t-t_1}{\tau}}\quad\text{i}\quad (V_{OH}-V_{OL})\left(e^{-\frac{t-t_2}{\tau}}-1\right).\]
Završna formula na slajdu izdvaja ostvareni porast tokom trajanja impulsa i njegovo kasnije eksponencijalno opadanje. Pri nenultom $V_{OL}$ polazni nivo i dalje ostaje deo rezultata.
